\documentclass[popl.tex]{subfiles}
\begin{document}
\section{ChatGPT repair prompt}\label{sec:chatgpt-prompt}

For the ChatGPT baseline in Fig.~\ref{fig:len_dist_prec}, we queried the OpenAI Responses API with model \texttt{gpt-5.4-mini}, \texttt{max\_output\_tokens=1024}, \texttt{store=false}, no temperature override, and no reasoning parameter. Each request used a single user message containing the following prompt, with \texttt{\$brokenTokens} replaced by the tokenized broken snippet:

\begin{lstlisting}[basicstyle=\ttfamily\scriptsize,breaklines=true]
   You are repairing a Python syntax error in tokenized form.

   The input and output are whitespace-separated lexer tokens, not source code.
   Identifiers and literals are abstracted as NAME, NUMBER, and STRING.
   Structural tokens such as NEWLINE, INDENT, DEDENT, and ENDMARKER must remain tokens.

   Return exactly one repaired token sequence in the same whitespace-separated token format.
   Return the complete repaired sequence, including all unchanged tokens from the input.
   Do not return only the changed token, only the edit, or a patch.
   Make the smallest syntactic repair that is most likely to be the author's intended code.
   Do not return markdown, XML tags, explanations, alternatives, or source code.

   Broken token sequence:
   $brokenTokens
\end{lstlisting}

The response was interpreted as a single whitespace-separated token sequence after stripping code fences, optional \texttt{Repair:} prefixes, surrounding quotes, and redundant whitespace. Across 2,238 repair instances in the test set, this costed roughly \$1 USD to evaluate as of June 2026 pricing.
\ifSubfilesClassLoaded{\par\clearpage\bibliography{../bib/acmart}}{}
\end{document}
